%% Springer Nature journal template -- sn-jnl (Version 3.1, December 2024)
%% Reference style: Math and Physical Sciences Numbered
\documentclass[pdflatex,sn-mathphys-num]{sn-jnl}

%%%% Packages
\usepackage{graphicx}
\usepackage{multirow}
\usepackage{amsmath,amssymb,amsfonts}
\usepackage{amsthm}
\usepackage{booktabs}
\usepackage{xcolor}
\usepackage[title]{appendix}
\usepackage{textcomp}

%%%% Theorem styles
\theoremstyle{thmstyleone}
\newtheorem{theorem}{Theorem}
\newtheorem{proposition}[theorem]{Proposition}

\theoremstyle{thmstyletwo}
\newtheorem{example}{Example}
\newtheorem{remark}{Remark}

\theoremstyle{thmstylethree}
\newtheorem{definition}{Definition}

\raggedbottom

\begin{document}

\title[Feature-Driven Meta-Learning for Decomposition Strategy Selection]{%
  When Does Seasonal Decomposition Help?\\
  A Feature-Driven Meta-Learning Study for Time Series Forecasting}

\author*[1]{\fnm{Diogo} \sur{Costa}}\email{diogocosta1104@gmail.com}

\affil*[1]{%
  \orgdiv{Faculty of Engineering},
  \orgname{University of Porto},
  \orgaddress{%
    \street{Rua Dr.\ Roberto Frias},
    \city{Porto},
    \postcode{4200-465},
    \country{Portugal}}}

\abstract{%
% TODO: Write after experiments are complete.
}

\keywords{Time series forecasting, Seasonal decomposition, STL,
  Meta-learning, Feature engineering, AutoML, Global forecasting models}

\maketitle

% ---------------------------------------------------------------------------
\section{Introduction}\label{sec:intro}
% TODO
% ---------------------------------------------------------------------------

% ---------------------------------------------------------------------------
\section{Related Work}\label{sec:related}
% TODO
% ---------------------------------------------------------------------------

% ---------------------------------------------------------------------------
\section{Methodology}\label{sec:methodology}
% ---------------------------------------------------------------------------

\subsection{Problem Formulation}\label{sec:problem}

Let $\mathbf{y} = (y_1, y_2, \ldots, y_T)$ be a univariate time series with
seasonal period $m$ and forecast horizon $h$.
Given a set of forecasting methods $\mathcal{M}$ and decomposition strategies
$\mathcal{D}$, the question is whether observable properties of $\mathbf{y}$ can
predict which pair $(d, M) \in \mathcal{D} \times \mathcal{M}$ minimises forecast
error for that series.

To study this, we extract a six-dimensional feature vector $\phi(\mathbf{y})$,
partition series into tercile groups along each feature dimension,
train all $(d, M)$ combinations within each group, and compare performance across
groups. If decomposition is beneficial, this signal should appear as a systematic
performance difference between feature terciles.

\subsection{Datasets}\label{sec:data}

We use two established benchmarks from the M-Competition series:
the M3~\cite{makridakis2000m3} and M4~\cite{makridakis2020m4} datasets,
restricted to monthly and quarterly frequencies.
Table~\ref{tab:datasets} summarises the resulting collection.

\begin{table}[ht]
\caption{Dataset summary. Horizon and season length follow the competition
definitions; series counts reflect the representative sample used for
benchmarking (see Section~\ref{sec:sampling}).}\label{tab:datasets}
\begin{tabular}{llccc}
\toprule
\textbf{Dataset} & \textbf{Frequency} & \textbf{Season ($m$)} &
\textbf{Horizon ($h$)} & \textbf{Series} \\
\midrule
M3 Monthly    & Monthly   & 12 & 18 & 750 \\
M4 Monthly    & Monthly   & 12 & 18 & 750 \\
M3 Quarterly  & Quarterly & 4  & 8  & 750 \\
M4 Quarterly  & Quarterly & 4  & 8  & 750 \\
\midrule
\multicolumn{4}{l}{Total} & 3{,}000 \\
\botrule
\end{tabular}
\end{table}

\subsubsection{Representative Sampling.}\label{sec:sampling}

The M3 and M4 pools each contain several thousand series per frequency, so we
draw 750 series per source dataset using stratified quantile sampling.
Series are ranked by each feature $\phi_k$ and divided into $S = 30$ equal-count
strata; the 750-series quota is distributed across strata proportionally and filled
by draws without replacement.
This ensures the sample reproduces the marginal feature distribution of the full pool.
A fixed random seed makes the sampling deterministic.

\subsection{Time Series Feature Extraction}\label{sec:features}

Six features characterise each series.
Where noted, they operate on STL residuals (Section~\ref{sec:stl}) to
focus on intrinsic series behaviour rather than removable seasonal structure.
All features are computed from the training observations only, with no access
to the test period.

\paragraph{Non-normality ($\phi_1$).}
Residuals $r_t$ from the STL fit are standardised as
$z_t = (r_t - \tilde{r}) / (1.4826 \cdot \mathrm{MAD}(r) + \varepsilon)$,
where $\tilde{r}$ is the median.
The D'Agostino--Pearson omnibus test~\cite{dagostino1971} is applied to
$\{z_t\}$, and its $p$-value is converted to a log-scale score:
\begin{equation}
  \phi_1 = \min\!\bigl(-\log_{10}(p_1 + 10^{-12}),\; 12\bigr).
  \label{eq:phi1}
\end{equation}
High values indicate strong departure from Gaussian residuals.

\paragraph{Nonlinearity ($\phi_2$).}
We assess whether lagged history contains nonlinear predictive structure by
comparing a linear and a nonlinear autoregressive model.
Let $\ell = \min(m, 12)$ be the embedding dimension.
We construct lag-feature matrices from $\{(y_{t-\ell},\ldots,y_{t-1}, y_t)\}$
and split the resulting dataset at three chronological split points
(55\%, 70\%, and 85\% of the series length) to obtain multiple train/test folds.
On each fold we fit a Ridge regression~\cite{hoerl1970ridge} ($\alpha = 1$) and
a Random Forest~\cite{breiman2001random} (40 trees, max depth 4).
The feature is the mean relative improvement of the nonlinear model:
\begin{equation}
  \phi_2 = \min\!\left(
    \max\!\left(0,\;
    \frac{\overline{\mathrm{MAE}}_{\mathrm{linear}} -
           \overline{\mathrm{MAE}}_{\mathrm{RF}}}
         {\overline{\mathrm{MAE}}_{\mathrm{linear}} + 10^{-8}}
    \right),\; 2
  \right),
  \label{eq:phi2}
\end{equation}
where $\overline{\mathrm{MAE}}$ denotes the mean across folds.
A value of zero indicates the Random Forest offers no improvement over ridge
regression; this is structurally common in linear datasets, as discussed in
Section~\ref{sec:nonlin_edge}.

\paragraph{Spectral entropy ($\phi_3$).}
We compute the power spectral density of the demeaned series via a periodogram
and derive its normalised Shannon entropy:
\begin{equation}
  \phi_3 = -\frac{\displaystyle\sum_{k} p_k \ln p_k}{\ln K}
          \;\in [0, 1],
  \label{eq:phi3}
\end{equation}
where $p_k = \hat{S}(f_k) / \sum_j \hat{S}(f_j)$ are the normalised spectral
weights and $K$ is the number of non-zero frequencies.
Low values indicate a spectrum concentrated at a few dominant frequencies
(strong, regular periodicity); high values indicate a diffuse, unpredictable
spectrum~\cite{spectral_entropy_ref}.

\paragraph{Evolving seasonality ($\phi_4$).}
To detect non-stationary seasonal behaviour, we compute the seasonal strength
(Equation~\ref{eq:seasstrength}) in rolling windows of length
$\max(36, 3m)$ for monthly and $\max(16, 4m)$ for quarterly series,
advancing by $\lfloor m/2 \rfloor$ observations per step.
The seasonal strength in each window is
\begin{equation}
  F_s^{(w)} = \max\!\left(0,\; 1 - \frac{\mathrm{Var}(R^{(w)})}
                                         {\mathrm{Var}(S^{(w)} + R^{(w)})}\right),
  \label{eq:seasstrength}
\end{equation}
where $S^{(w)}$ and $R^{(w)}$ are the STL seasonal and residual components in
window $w$.
The feature is the standard deviation of the sequence $\{F_s^{(w)}\}$:
a high value indicates that seasonal strength changes substantially over time.

\paragraph{Structural break strength ($\phi_5$).}
We detect level shifts in the nonseasonal component $T_t + R_t$ using the
PELT algorithm~\cite{killick2012} with an RBF cost function and a
$\log(n)$-proportional penalty.
Let $\mathcal{B}$ be the set of detected breakpoints in the central 60\% of the
series (to exclude boundary effects).
For each candidate break $b \in \mathcal{B}$, the mean-shift contrast is
\begin{equation}
  \Delta_b = \frac{|\bar{x}_{<b} - \bar{x}_{\geq b}|}
                  {1.4826 \cdot \mathrm{MAD}(T + R) + \varepsilon},
  \label{eq:breakcontrast}
\end{equation}
and $\phi_5 = \min(\max_b \Delta_b,\; 20)$.
This measures the magnitude of the largest structural shift relative to
the scale of the nonseasonal component.

\paragraph{Conditional heteroscedasticity ($\phi_6$).}
We test for ARCH effects~\cite{engle1982arch} in the normalised STL residuals
$\{z_t\}$ using Engle's LM test with lag order $\min(12, \lfloor n/5 \rfloor)$
and return
\begin{equation}
  \phi_6 = \min\!\bigl(-\log_{10}(p_6 + 10^{-12}),\; 12\bigr),
  \label{eq:phi6}
\end{equation}
where $p_6$ is the LM test $p$-value.
High values indicate serially correlated residual variance (volatility
clustering).

\subsubsection{Feature Bucket Assignment.}\label{sec:buckets}

For each feature $\phi_k$ and frequency, all selected series are jointly ranked
by their feature value and assigned to tercile buckets: \textit{Low} (bottom
third), \textit{Medium} (middle third), and \textit{High} (top third).
The bucket boundaries are the $33.3^{\text{rd}}$ and $66.7^{\text{th}}$
percentiles computed across both M3 and M4 series combined.
This partitioning drives the feature-conditioned finetuning strategy described
in Section~\ref{sec:finetuning}.

\subsection{STL Decomposition Framework}\label{sec:stl}

Seasonal-Trend decomposition using Loess (STL)~\cite{cleveland1990stl}
factorises a time series as
\begin{equation}
  y_t = T_t + S_t + R_t,
  \label{eq:stl}
\end{equation}
where $T_t$ is a locally smooth trend, $S_t$ is a periodic seasonal component
(constrained to be approximately zero-mean within each period), and $R_t$ is
the remainder.
We use the robust variant of STL (down-weighting outlier observations during
the Loess smoothing passes), implemented via \texttt{statsmodels}~\cite{seabold2010}.

We evaluate three strategies for exploiting this decomposition:

\begin{description}
  \item[\textbf{Direct} (\texttt{without\_stl}).] The forecasting model operates
    directly on the raw series $\mathbf{y}$. No decomposition is applied.

  \item[\textbf{STL--Seasonal-Naive} (\texttt{stl\_seasonal\_naive}).] The model
    is trained on the nonseasonal component $N_t = T_t + R_t$, producing a
    forecast $\hat{N}_{t+1:t+h}$. The seasonal component is extrapolated by
    repeating the last observed full seasonal cycle:
    $\hat{S}_{t+j} = S_{t+j-m}$ for $j = 1, \ldots, h$.
    The final forecast is $\hat{y}_{t+j} = \hat{N}_{t+j} + \hat{S}_{t+j}$.

  \item[\textbf{STL--All-Components} (\texttt{stl\_model\_all\_components}).] Three
    separate model instances are trained on $T_t$, $S_t$, and $R_t$ respectively,
    yielding component forecasts $\hat{T}$, $\hat{S}$, and $\hat{R}$.
    The final forecast is the sum:
    \begin{equation}
      \hat{y}_{t+j} = \hat{T}_{t+j} + \hat{S}_{t+j} + \hat{R}_{t+j},
      \quad j = 1, \ldots, h.
      \label{eq:recompose}
    \end{equation}
\end{description}

\subsection{Forecasting Model Families}\label{sec:models}

We benchmark four families across two training paradigms.

\paragraph{Statistical (\textit{per-series}).}
AutoARIMA~\cite{hyndman2008automatic}, AutoSARIMA (forced $D=1$), and
AutoETS~\cite{hyndman2008exponential}, all implemented via
\texttt{statsforecast}~\cite{olivares2022statsforecast}.

\paragraph{Machine learning (\textit{global}).}
AutoRidge, AutoLinearRegression, AutoXGBoost~\cite{chen2016xgboost}, and
AutoRandomForest, trained jointly on all cohort series via sliding-window
lag features (\texttt{MLForecast}~\cite{garza2022mlforecast}, Optuna
HPO~\cite{akiba2019optuna}, 100 trials).

\paragraph{Neural (\textit{global}).}
AutoNLinear, AutoNHITS~\cite{challu2023nhits}, and AutoLSTM
(\texttt{NeuralForecast}~\cite{olivares2023neuralforecast}).

\paragraph{Transformers (\textit{global}).}
AutoTFT~\cite{lim2021tft} and AutoPatchTST~\cite{nie2022patchtst}
(\texttt{NeuralForecast}).

\subsection{Bucket-Specialised Training}\label{sec:finetuning}

For global model families (ML, Neural, Transformers), we investigate whether
\emph{regime-specialised} global models outperform a single global model trained
on all 750 series.
The hypothesis is that series sharing the same feature-tercile level exhibit
sufficiently homogeneous dynamics to benefit from a dedicated model.

For each task configuration (feature $\times$ frequency $\times$ decomposition
$\times$ model), three independent global models are trained---one per tercile
bucket (\textit{Low}, \textit{Medium}, \textit{High})---each using
\emph{exclusively} the $\approx\!250$ series in its bucket, and evaluated
\emph{exclusively} on the test series from the same bucket.
Performance is measured as the finetuning delta
$\Delta_{\mathrm{FT}} = \mathrm{RelNaive}^{(10)}_{\mathrm{global}} -
\mathrm{RelNaive}^{(10)}_{\mathrm{bucket}}$ (positive = specialised model wins).
The global model (all 750 series, no bucket restriction) serves as the
comparison baseline throughout.

Statistical models are excluded as they fit independently per series and are
unaffected by the training-pool composition.
For global families, each (decomposition, model, feature) configuration is
evaluated twice: once with all 750 series as the global baseline, and once per
tercile bucket under the specialised condition.

\subsection{Evaluation Protocol}\label{sec:evaluation}

\subsubsection{Cross-Validation.}
We use rolling-origin cross-validation~\cite{tashman2000} with $W = 3$ windows
per series, advancing by $h$ observations per window.
For each series $i$ and window $w$, we define a training set ending at cutoff
$c_i^{(w)}$ and a test set of exactly $h$ future observations.

\subsubsection{Metric.}
The primary metric is the \emph{Relative Naive} (RelNaive), defined as
\begin{equation}
  \mathrm{RelNaive} =
    \frac{\mathrm{MAE}(\mathbf{y},\,\hat{\mathbf{y}})}
         {\mathrm{MAE}(\mathbf{y},\,\hat{\mathbf{y}}^{\mathrm{naive}})},
  \label{eq:relnai}
\end{equation}
where $\hat{\mathbf{y}}^{\mathrm{naive}}$ is the seasonal naive forecast
($\hat{y}_{t+j}^{\mathrm{naive}} = y_{t+j-m}$, $j = 1, \ldots, h$).
A value below 1 indicates the model outperforms the seasonal naive baseline.
We report both unclipped ($\mathrm{RelNaive}$) and clipped
($\mathrm{RelNaive}^{(10)}$, capped at 10) variants; the latter reduces the
influence of catastrophic failures on aggregate rankings.
Aggregate statistics are computed as the mean over all series and windows within
each reporting group (feature, decomposition method, model family, training scope).

The \emph{STL delta} $\Delta_{\mathrm{STL}}$ measures the improvement from
applying STL decomposition relative to the Direct strategy:
\begin{equation}
  \Delta_{\mathrm{STL}} =
    \mathrm{RelNaive}^{(10)}_{\mathrm{Direct}} -
    \mathrm{RelNaive}^{(10)}_{\mathrm{STL}},
  \label{eq:stl_delta}
\end{equation}
so a positive value indicates that STL decomposition reduces error.
The \emph{specialisation delta} $\Delta_{\mathrm{FT}}$ measures the improvement
from training a bucket-specialised model over the global baseline:
\begin{equation}
  \Delta_{\mathrm{FT}} =
    \mathrm{RelNaive}^{(10)}_{\mathrm{global}} -
    \mathrm{RelNaive}^{(10)}_{\mathrm{bucket}},
  \label{eq:ft_delta}
\end{equation}
with a positive value indicating that regime-specialised training outperforms
the global model on the same series.

Statistical significance of aggregate deltas is assessed with the
Wilcoxon signed-rank test~\cite{wilcoxon1945} applied at the series level,
with $p$-values corrected for multiple comparisons using the Holm--Bonferroni
method~\cite{holm1979}.

% ---------------------------------------------------------------------------
\section{Results}\label{sec:results}
% ---------------------------------------------------------------------------

% ---------------------------------------------------------------------------
\subsection{Global Baseline Performance}\label{sec:baseline}
% ---------------------------------------------------------------------------

All 3{,}000 series (1{,}500 monthly, 1{,}500 quarterly) are evaluated over three
rolling cross-validation windows.
Table~\ref{tab:top10} ranks the global model configurations by mean
$\mathrm{RelNaive}^{(10)}$ (all series, no bucket restriction).
\textbf{AutoETS with Direct (no decomposition)} is the best single configuration
(mean $0.888$, median $0.791$), and the only one to beat the na\"ive forecast at
the mean level without preprocessing.
The next nine are dominated by STL-SN paired with statistical or transformer
models.

\begin{table}[ht]
\caption{Top-10 global model configurations by mean $\mathrm{RelNaive}^{(10)}$
(lower is better).
Direct = no decomposition; STL-SN = STL--Seasonal-Na\"ive;
STL-AC = STL--All-Components.}\label{tab:top10}
\begin{tabular*}{\textwidth}{@{\extracolsep\fill}llllcc}
\toprule
Rank & Family & Model & Decomp.\ & Mean & Median \\
\midrule
1  & Statistical  & AutoETS               & Direct  & 0.888 & 0.791 \\
2  & Statistical  & AutoARIMA             & STL-SN  & 0.930 & 0.887 \\
3  & Neural       & AutoNLinear           & STL-SN  & 0.931 & 0.889 \\
4  & Statistical  & AutoSARIMA            & STL-SN  & 0.931 & 0.889 \\
5  & ML           & AutoLinearRegression  & STL-SN  & 0.933 & 0.893 \\
6  & ML           & AutoRidge             & STL-SN  & 0.933 & 0.889 \\
7  & Transformer  & AutoPatchTST          & STL-SN  & 0.933 & 0.891 \\
8  & Statistical  & AutoETS               & STL-SN  & 0.935 & 0.918 \\
9  & Transformer  & AutoPatchTST          & STL-AC  & 0.936 & 0.826 \\
10 & Transformer  & AutoTFT               & STL-SN  & 0.936 & 0.919 \\
\botrule
\end{tabular*}
\end{table}

A persistent divergence between mean and median runs throughout the study:
under STL decomposition the ML and Neural families have mean
$\mathrm{RelNaive}^{(10)} > 1$ (1.12--1.19) while their medians remain near or
below 1.
A heavy right tail of catastrophic failures on a minority of series pulls the
mean above the na\"ive baseline even when most series benefit.
Both statistics are therefore reported throughout.

% ---------------------------------------------------------------------------
\subsection{Effect of STL Decomposition}\label{sec:stl_results}
% ---------------------------------------------------------------------------

Table~\ref{tab:stl_delta} reports $\Delta_{\mathrm{STL}}$ for the global model.
Positive values indicate STL reduces error relative to Direct.

\begin{table}[ht]
\caption{STL delta $\Delta_{\mathrm{STL}}$ by decomposition and model family
(global model; mean and median over all series and windows;
win rate = fraction of series where STL helps).}\label{tab:stl_delta}
\begin{tabular*}{\textwidth}{@{\extracolsep\fill}llccc}
\toprule
Decomposition & Family & Mean $\Delta$ & Median $\Delta$ & Win Rate \\
\midrule
STL-AC & Transformer  & $+0.039$ & $+0.141$ & $68.5\%$ \\
STL-SN & Transformer  & $+0.063$ & $+0.091$ & $67.9\%$ \\
STL-AC & Neural       & $-0.119$ & $+0.063$ & $57.6\%$ \\
STL-SN & Neural       & $-0.109$ & $+0.027$ & $55.4\%$ \\
STL-AC & ML           & $-0.191$ & $+0.023$ & $52.5\%$ \\
STL-SN & ML           & $-0.194$ & $-0.002$ & $49.5\%$ \\
STL-AC & Statistical  & $-0.028$ & $-0.085$ & $39.8\%$ \\
STL-SN & Statistical  & $-0.009$ & $-0.109$ & $37.1\%$ \\
\botrule
\end{tabular*}
\end{table}

\paragraph{Transformers} benefit robustly from STL decomposition: both
STL-SN and STL-AC yield positive mean \emph{and} median deltas, with win
rates near 68\%.
The improvement is strongest under STL-AC (median $+0.141$), suggesting that
training separate global submodels for trend, seasonal, and residual components
aligns well with PatchTST and TFT inductive biases.

\paragraph{Neural and ML models} show the mean--median divergence at its
most pronounced.
For Neural, the median is positive (STL-AC: $+0.063$; STL-SN: $+0.027$) but
the mean is deeply negative ($-0.119$, $-0.109$): STL helps most series but
fails badly on those with non-stationary seasonal patterns where the
decomposition itself is unreliable.
ML is worse: STL-SN mean $-0.194$, median essentially zero ($-0.002$).

\paragraph{Statistical models} are consistently \emph{harmed} by both
decomposition strategies (win rates $\approx 38$--$40\%$, negative mean and
median).
AutoARIMA, AutoSARIMA, and AutoETS already model trend and seasonality
internally; applying STL preprocessing introduces redundant transformation and
additional parameter uncertainty, degrading rather than improving performance.

Frequency effects are marginal: quarterly series show slightly larger median
gains (STL-AC median $+0.024$ vs.\ $+0.033$ monthly) with no directional
difference.

% ---------------------------------------------------------------------------
\subsection{Bucket-Specific Finetuning}\label{sec:ft_results}
% ---------------------------------------------------------------------------

\subsubsection{Overall Results}\label{sec:ft_overall}

Across 1{,}449{,}738 series-window observations, the specialised models achieve
median $\Delta_{\mathrm{FT}} = +0.050$, mean $= -0.029$, win rate $= 56.1\%$.
The mean--median gap mirrors the pattern seen for STL: bucket specialisation
benefits the majority but fails badly on a tail of series, dragging the mean
negative.

Medium-bucket series show the most consistent improvement (median $+0.059$,
win rate $56.8\%$); Low and High are comparable (median $+0.048$ and $+0.046$).
The Medium advantage reflects that these are the modal-regime series---the
specialised model is not pulled toward the extremes present in either tail.

\subsubsection{Non-Linear vs.\ Linear Models: A Critical Split}\label{sec:linvnonlin}

The clearest split is between non-linear and linear models:

\begin{table}[ht]
\caption{Bucket-specific finetuning delta by model type (aggregated over all
families, buckets, decompositions, and series).}\label{tab:lin_nonlin}
\begin{tabular}{lccc}
\toprule
Model type & Mean $\Delta_{\mathrm{FT}}$ & Median $\Delta_{\mathrm{FT}}$ & Win rate \\
\midrule
Non-linear & $+0.082$ & $+0.092$ & $60.1\%$ \\
Linear     & $-0.251$ & $-0.014$ & $48.1\%$ \\
\botrule
\end{tabular}
\end{table}

\noindent
\textbf{Non-linear models benefit} because focusing training on $\approx\!250$
regime-homogeneous series improves generalisation within that regime without
sacrificing coverage of the full distribution.
Both mean and median are positive, so the gain is robust.

\noindent
\textbf{Linear models are harmed}: $\approx\!250$ training points are
insufficient for stable coefficient estimation, and the bucket-induced
distribution shift worsens the problem.
The sub-50\% win rate ($48.1\%$) confirms a directional harm, not noise;
Section~\ref{sec:significance} verifies significance.

Table~\ref{tab:per_model} gives the per-model breakdown.

\begin{table}[ht]
\caption{Bucket-specific finetuning delta per model (aggregated over all
buckets, decompositions, features, and frequencies).
Models are classified as linear (L) or non-linear (NL).}\label{tab:per_model}
\begin{tabular*}{\textwidth}{@{\extracolsep\fill}lllccc}
\toprule
Family & Model & Type & Mean $\Delta$ & Median $\Delta$ & Win rate \\
\midrule
ML          & AutoXGBoost           & NL & $+0.236$ & $+0.187$ & $65.3\%$ \\
Neural      & AutoLSTM              & NL & $+0.231$ & $+0.162$ & $64.2\%$ \\
ML          & AutoRandomForest      & NL & $+0.146$ & $+0.155$ & $61.6\%$ \\
Neural      & AutoNHITS             & NL & $-0.031$ & $+0.076$ & $59.2\%$ \\
Transformer & AutoTFT               & NL & $-0.076$ & $+0.059$ & $57.5\%$ \\
Transformer & AutoPatchTST          & NL & $-0.013$ & $+0.014$ & $52.6\%$ \\
\midrule
ML          & AutoRidge             & L  & $-0.327$ & $-0.004$ & $49.5\%$ \\
ML          & AutoLinearRegression  & L  & $-0.327$ & $-0.002$ & $49.7\%$ \\
Neural      & AutoNLinear           & L  & $-0.099$ & $-0.029$ & $45.1\%$ \\
\botrule
\end{tabular*}
\end{table}

\noindent
AutoXGBoost and AutoLSTM show the largest gains (median $+0.187$,
win $65.3\%$; median $+0.162$, win $64.2\%$ respectively), benefiting from
the focused within-bucket training distribution.
AutoNHITS and the transformers have positive medians but negative means,
indicating fat-tailed distributions.
AutoRidge and AutoLinearRegression suffer the largest harms (mean $\Delta =
-0.327$), a direct consequence of coefficient instability on small training
pools.
\textbf{Bucket specialisation should be restricted to non-linear models.}

\subsubsection{Interaction with STL Decomposition}\label{sec:ft_stl_interact}

Finetuning interacts differently with each decomposition strategy:

\begin{table}[ht]
\caption{Bucket finetuning delta by decomposition strategy (all non-statistical
model families and buckets combined).}\label{tab:ft_decomp}
\begin{tabular}{lccc}
\toprule
Decomposition & Mean $\Delta_{\mathrm{FT}}$ & Median $\Delta_{\mathrm{FT}}$ & Win rate \\
\midrule
Direct   & $-0.156$ & $+0.108$ & $60.4\%$ \\
STL-SN   & $+0.044$ & $+0.032$ & $55.6\%$ \\
STL-AC   & $+0.026$ & $+0.021$ & $52.2\%$ \\
\botrule
\end{tabular}
\end{table}

\noindent
Direct gives the highest median ($+0.108$) but the worst mean ($-0.156$):
without seasonal preprocessing, a model trained on $\approx\!250$ series
overfits to the idiosyncratic structure of its bucket.
\textbf{STL-SN is the most robust pairing} (mean $+0.044$, median $+0.032$):
removing seasonality first produces a more stationary target, reducing
catastrophic failures.

% ---------------------------------------------------------------------------
\subsection{The Feature Nonlinearity Edge Case}\label{sec:nonlin_edge}
% ---------------------------------------------------------------------------

The \texttt{feature\_nonlinearity} ($\phi_2$) feature exposes a structural
limitation of tercile-based bucketing.
Because $\phi_2$ is zero-floored (Equation~\ref{eq:phi2}), the M3/M4 series
accumulate at exactly zero: 81.8\% monthly, 92.8\% quarterly.
Both tercile boundaries therefore coincide at $q_{1/3} = q_{2/3} = 0$, so all
zero-valued series fall into Low and all non-zero series into High, leaving the
Medium bucket empty (Table~\ref{tab:nonlin_buckets}).

\begin{table}[ht]
\caption{Bucket sizes for \texttt{feature\_nonlinearity} ($\phi_2$).
$q_{1/3} = q_{2/3} = 0$ forces binary Low/High stratification.}\label{tab:nonlin_buckets}
\begin{tabular}{lcccc}
\toprule
Frequency & Low & Medium & High & Total \\
\midrule
Monthly   & 1{,}227 (81.8\%) & \textbf{0} & 273 (18.2\%) & 1{,}500 \\
Quarterly & 1{,}392 (92.8\%) & \textbf{0} & 108 (7.2\%)  & 1{,}500 \\
\botrule
\end{tabular}
\end{table}

The 27 affected task configurations (1.9\% of all bucket tasks) are excluded
from comparative analyses.
The binary split retains a meaningful interpretation: Low = series where
lagged ridge regression is near-optimal; High = the minority with
genuinely exploitable nonlinear structure.
This is in fact the sharpest regime contrast among all six features.

% ---------------------------------------------------------------------------
\subsection{Statistical Significance}\label{sec:significance}
% ---------------------------------------------------------------------------

Wilcoxon signed-rank tests with Holm--Bonferroni correction were applied to
each (model, decomposition, bucket) combination ($n = 81$ tests total).
Results:

\begin{itemize}
  \item \textbf{Significantly helps:} 55/81 (67.9\%),
        all $p_{\mathrm{Holm}} < 0.001$.
  \item \textbf{Significantly harms:} 23/81 (28.4\%),
        all $p_{\mathrm{Holm}} < 0.001$.
  \item \textbf{Neutral (not significant):} 3/81 (3.7\%).
\end{itemize}

The five most improved combinations are all non-linear models paired with
the Medium bucket and STL decomposition (Table~\ref{tab:top5_sig}).
The STL-SN + Medium + AutoXGBoost combination is the single strongest result
in the study (median $\Delta_{\mathrm{FT}} = +0.334$, win rate $72.3\%$).

\begin{table}[ht]
\caption{Top-5 significantly improved model $\times$ bucket combinations
(Wilcoxon + Holm, all $p_{\mathrm{Holm}} < 0.001$).}\label{tab:top5_sig}
\begin{tabular*}{\textwidth}{@{\extracolsep\fill}llllcc}
\toprule
Family & Model & Decomp.\ & Bucket & Med.\ $\Delta$ & Win rate \\
\midrule
ML          & AutoXGBoost      & STL-SN & Medium & $+0.334$ & $72.3\%$ \\
ML          & AutoXGBoost      & STL-AC & Medium & $+0.308$ & $66.9\%$ \\
Neural      & AutoLSTM         & STL-SN & Medium & $+0.299$ & $69.2\%$ \\
ML          & AutoRandomForest & STL-AC & Medium & $+0.291$ & $64.4\%$ \\
ML          & AutoRandomForest & STL-SN & Medium & $+0.290$ & $68.2\%$ \\
\botrule
\end{tabular*}
\end{table}

The significantly harmed combinations are exclusively linear models and
AutoNLinear, with the worst harm for AutoNLinear under STL-AC
(median $\Delta_{\mathrm{FT}} = -0.088$, win rate $34.5\%$).
Linear models under STL-AC suffer doubly: the global model must simultaneously
fit trend, seasonal, and residual components of $\approx\!250$ series, giving
each linear submodel too few degrees of freedom while compounding errors
through recomposition.

% ---------------------------------------------------------------------------
\subsection{Summary of Findings}\label{sec:summary}
% ---------------------------------------------------------------------------

\begin{enumerate}
  \item \textbf{STL decomposition effects are family-specific}
    (Section~\ref{sec:stl_results}).
    Transformers benefit robustly ($\Delta_{\mathrm{STL}}$ positive in mean
    \emph{and} median, win rate $\approx 68\%$).
    Neural and ML models show positive medians but severely negative means due
    to STL failure on non-stationary series.
    Statistical models are uniformly harmed---their internal trend-seasonal
    modelling makes STL preprocessing redundant.

  \item \textbf{Bucket specialisation improves non-linear models, harms linear
    ones} (Section~\ref{sec:linvnonlin}).
    Non-linear architectures gain median $+0.092$ (win rate 60.1\%); linear
    models lose median $-0.014$ (win rate 48.1\%).
    The split is statistically significant in 67.9\% of tested configurations.

  \item \textbf{The Medium bucket is most reliable; STL-SN is the safest
    pairing} (Sections~\ref{sec:ft_overall}, \ref{sec:ft_stl_interact}).
    Medium-bucket series yield the most consistent specialisation gains.
    STL-SN before bucket training gives positive mean \emph{and} median delta,
    eliminating the catastrophic failures seen under Direct.

  \item \textbf{Feature nonlinearity collapses to binary bucketing on M-competition
    data} (Section~\ref{sec:nonlin_edge}).
    The zero-floor of $\phi_2$ causes $q_{1/3} = q_{2/3} = 0$, emptying the
    Medium bucket and reducing stratification to Low/High for 1.9\% of tasks.
    Tercile-based partitioning requires a minimum-support guard when features
    have degenerate distributions.
\end{enumerate}

% ---------------------------------------------------------------------------
\section{Discussion}\label{sec:discussion}
% TODO
% ---------------------------------------------------------------------------

% ---------------------------------------------------------------------------
\section{Conclusion}\label{sec:conclusion}
% TODO
% ---------------------------------------------------------------------------

\backmatter

\bmhead{Acknowledgements}
% TODO

\section*{Declarations}

\begin{itemize}
  \item \textbf{Funding:} Not applicable.
  \item \textbf{Conflict of interest:} The author declares no competing interests.
  \item \textbf{Ethics approval:} Not applicable.
  \item \textbf{Data availability:} The M3 and M4 datasets are publicly available
    at \url{https://forecasters.org/resources/time-series-data/}.
  \item \textbf{Code availability:} Source code will be made available upon
    acceptance.
  \item \textbf{Author contribution:} D.C. conceived and designed the study,
    implemented the pipeline, ran all experiments, and wrote the manuscript.
\end{itemize}

\bibliography{references}

\end{document}
